\par\Needspace{6\baselineskip}\medskip
\begin{center}\hypertarget{native-177-article130}{}\pdfbookmark[2]{130.}{article-130}130.\end{center}
Die zehn ab Artikel 124 erläuterten Verfahren setzen voraus, dass bereits Näherungswerte für
die Abstände des Himmelskörpers von der Erde oder für die Lage der Bahnebene vorliegen. Diese
Voraussetzung bereitet offenbar keine Schwierigkeiten, wenn die bereits näherungsweise bekannten
Abmessungen der Bahn anhand weiter auseinanderliegender Beobachtungen korrigiert werden sollen. Die
Näherungswerte können etwa aus einer früheren, auf anderen Beobachtungen beruhenden Rechnung
stammen. Damit ist noch offen, wie man die erste Rechnung beginnen kann, wenn alle Abmessungen der
Bahn noch völlig unbekannt sind. Das ist der bei weitem wichtigste und schwierigste Fall unserer
Aufgabe, wie schon die entsprechende Aufgabe in der Theorie der Kometen vermuten lässt. Es ist
hinlänglich bekannt, wie lange jene Aufgabe die Mathematiker beschäftigte und wie viele vergebliche
Versuche sie veranlasste. Soll unsere Aufgabe als richtig gelöst gelten, so muss eine Lösung nach
dem ab Artikel 119 erläuterten Vorgehen offenbar die folgenden Bedingungen erfüllen. \emph{Erstens}
sind die Größen $x,y$ so zu wählen, dass sich ihre Näherungswerte aus der Natur der Aufgabe selbst
gewinnen lassen, zumindest solange die heliozentrische Bewegung des Himmelskörpers zwischen den
Beobachtungen nicht zu groß ist. \emph{Zweitens} dürfen kleine Änderungen der Größen $x,y$ in den
daraus abgeleiteten Größen keine zu großen Änderungen hervorrufen. Sonst könnten etwaige Fehler in
den angenommenen Werten der ersteren verhindern, dass man die letzteren noch als Näherungswerte
ansehen kann. \emph{Drittens} schließlich verlangen wir, dass die Rechenschritte, die von den Größen
$x,y$ nach und nach zu $X,Y$ führen, nicht zu langwierig werden.

Diese Bedingungen liefern einen Maßstab, nach dem sich die Güte eines jeden Verfahrens beurteilen
lässt. Bei häufiger Anwendung zeigt sie sich noch deutlicher. Das Verfahren, das wir nun erläutern
wollen und das in gewisser Weise als der wichtigste Teil dieses Werks gelten kann, erfüllt diese
Bedingungen so gut, dass es anscheinend nichts weiter \SourceLeaf{178}zu wünschen übrig lässt. Bevor wir es in
der für die Praxis geeignetsten Form erläutern, schicken wir einige Überlegungen voraus. Sie sollen
gleichsam den Zugang zu diesem Verfahren erhellen und eröffnen, der sonst vielleicht weniger klar
und weniger leicht zu finden wäre.
